% ============================================================
%  从Transformer数学结构到Agent工程约束：一个收敛理论框架
%  arXiv 初稿
% ============================================================

\documentclass[11pt,a4paper]{ctexart}

% ---- 页面设置 ----
\usepackage[left=3cm,right=3cm,top=2.5cm,bottom=2.5cm]{geometry}

% ---- 数学与符号 ----
\usepackage{amsmath,amssymb,amsthm}
\usepackage{mathtools}

% ---- 图形与表格 ----
\usepackage{graphicx}
\usepackage{booktabs}
\usepackage{array}
\usepackage{caption}
\usepackage{subcaption}
\usepackage{tikz}
\usetikzlibrary{arrows,shapes,positioning,fit,backgrounds}

% ---- 引用与链接 ----
\usepackage[colorlinks=true,linkcolor=blue,citecolor=blue,urlcolor=blue]{hyperref}
\usepackage{url}

% ---- 定理环境 ----
\newtheorem{theorem}{定理}[section]
\newtheorem{proposition}[theorem]{命题}
\newtheorem{lemma}[theorem]{引理}
\newtheorem{corollary}[theorem]{推论}
\newtheorem{definition}[theorem]{定义}
\newtheorem{remark}[theorem]{说明}

% ---- 其他 ----
\usepackage{enumitem}
\usepackage{fancyhdr}
\pagestyle{fancy}
\fancyhf{}
\fancyhead[L]{\small 从Transformer数学结构到Agent工程约束}
\fancyhead[R]{\small 初稿~~\today}
\fancyfoot[C]{\thepage}
\renewcommand{\headrulewidth}{0.4pt}

% ---- 标题 ----
\title{\textbf{从Transformer数学结构到Agent工程约束\\[4pt]
        ——一个收敛理论框架}}
\author{}
\date{\today}

\begin{document}

\maketitle

\begin{center}
\rule{0.8\textwidth}{0.5pt}
\end{center}

% ============================================================
\begin{abstract}
\noindent
本文论证：大语言模型Agent的一系列已知限制——任务分解的结构性脆弱、记忆缺乏时间动态厚度、全自动闭环的不可达——不是工程上待解决的暂时问题，而是Transformer架构的数学结构所决定的必然结果。

论证从Transformer注意力的数学形式出发：(1) 将注意力机制表述为语义空间中的非参数贝叶斯推断，引入米田嵌入刻画该语义空间的完备结构；(2) 基于参数空间有效维度的经验证据，论证模型学到的仅是完整语义空间的截断投影；(3) 提出训练收敛与测试收敛的双模分类，统一解释模型能力与Agent行为的来源与边界；(4) 从空间结构与顺序结构的异质性出发，论证LLM记忆中不存在长短期之间的连续过渡；(5) 运用三段论结构推导上下文驱动的任务分解不可靠是结构必然；(6) 从收敛边界推导全自动Agent闭环的原则不可闭合性——根需求必须来自系统外部。

本文在相关工作部分逐一标注了与已有文献的差异化路径，在讨论部分列出了可检验的推论及其证伪条件，在局限部分诚实标注了当前论证的薄弱环节。本文的目标不是断言该框架已完备，而是提出一个足以组织现有经验观察的统一理论结构，供社区检验、修正或反驳。
\end{abstract}

% ============================================================

\section{引言}

\subsection{问题：Agent的经验观察缺乏统一理论解释}

大语言模型（LLM）在作为Agent运行时，暴露出一系列系统性的、可复现的失败模式：

\begin{itemize}[leftmargin=*]
    \item \textbf{任务分解不可靠。} 模型在多步规划中系统性偏离目标，早期全自动Agent（如AutoGPT）几乎全部在开放域长程任务中失败；即便是配备了成熟脚手架的生产级Agent（如Cursor、Claude Code），仍需在关键决策点设置人工确认。
    \item \textbf{上下文窗口的"迷失中间"效应。} 窗口中间位置的信息检索准确率系统性低于首尾——这不是工程实现不优，而是softmax注意力的结构特征。
    \item \textbf{记忆缺乏时间深度。} 模型在长对话中混淆时序、遗忘早期决策、无法区分"三小时前说的"和"刚才说的"——这些失败无法仅靠扩大上下文窗口消除。
    \item \textbf{全自动闭环失败。} 从目标设定到任务完成的全自动闭环，在开放域中从未被可靠实现；而每次"近乎成功"的运行都依赖于人类在过程中隐性注入意图澄清和需求仲裁。
\end{itemize}

现有文献将这些失败归类为工程问题——提出更好的分解策略\cite{LEAD2025}、更长的上下文窗口、更优的记忆模块\cite{CognitiveMemory2025}、更复杂的多Agent协作框架\cite{MACI2025}。这些工作在各自的设定下取得了经验改进。

\textbf{本文的立场与此不同。} 我们认为，上述现象不是待解决的工程问题——它们是架构数学结构决定的边界。一个系统在既有架构内的最优工程方案，无法越过其数学结构划定的边界。不理解这些边界的来源和性质，就无法在正确的范围内设计系统。

\subsection{本文的论证路径}

本文的论证结构如下：

\begin{enumerate}[leftmargin=*,label=§\arabic*]
    \item \textbf{理论基础}（§2--§3）：从Transformer注意力的数学形式出发，引入范畴论（米田引理）和贝叶斯推断作为理论工具，将学习过程表述为函子，建立"收敛"概念的形式基础。
    \item \textbf{记忆的结构约束}（§4）：从参数空间与上下文空间的数学异质性出发，推导记忆的结构性二分，论证不存在长短期之间的连续过渡。
    \item \textbf{测试收敛理论}（§5）：提出Agent行为的真正支配原理——测试收敛——并分析其与训练收敛的根本差异，包括保证性缺失和有效性条件。
    \item \textbf{上下文边界与任务分解}（§6--§8）：从上下文窗口的信息边界出发，推导任务分解的结构不可靠性、层级收敛的性质、以及收敛的终止条件。
    \item \textbf{全自动闭环的不可闭合性}（§9）：给出不可闭合的精确定义和范围限定，区分目标设定/任务分解/执行闭环三层，并从收敛空间有界性给出论证。
    \item \textbf{工程含义}（§10）：导出Agent工程设计的五项根本原则，明确当前架构下的可行域边界。
\end{enumerate}

\subsection{各节论证强度分层说明}

\textbf{重要：本文各节处于不同的论证层次。} 建议读者根据不同层次的论证强度，选择性地评估和采纳本文的各个部分。

\begin{itemize}[leftmargin=*]
    \item \textbf{类比层次。} §2中米田嵌入与贝叶斯推断的统一目前处于类比层面——两者在结构上的对应关系已被识别，但尚未给出严格的数学对偶构造。§12.3讨论了通向形式化的开放问题。该层提供深层数学直觉，但不承载本文的核心论证链。

    \item \textbf{形式论证层次。} §3--§9的核心论证链——训练收敛与测试收敛的区分、记忆的结构性二分、上下文驱动的分解决策不可靠、全自动闭环的不可闭合——依赖两个有独立经验证据的前提：(a) 参数空间有效维度远小于名义维度\cite{Li2018IntrinsicDim,Aghajanyan2020}；(b) 上下文窗口中的信息不经过训练收敛的筛选。这两个前提均有独立于本文框架的经验支持。读者可选择在不接受§2类比构造的前提下独立评估§3--§9。

    \item \textbf{经验推论层次。} §10的工程原则和§12.1的可检验推论是从理论框架导出的经验命题，其有效性可独立于理论框架进行检验。如果推论被实验证伪，框架的部分或全部需要修正。
\end{itemize}

\subsection{贡献}

本文将以下五个贡献组织于一个统一的论证结构中：

\begin{enumerate}[leftmargin=*,label=\textbf{C\arabic*}.]
    \item 将注意力机制表述为米田嵌入语义空间中的贝叶斯推断，为分布语义学提供范畴论基础框架。
    \item 提出训练收敛与测试收敛的双模分类，统一解释模型能力（训练收敛保证）与Agent行为（测试收敛驱动、缺乏收敛保证）的来源与边界。
    \item 从结构必然性角度论证LLM记忆中不存在长短期之间的连续过渡——记忆只有参数空间的静态存储和上下文空间的顺序动态，没有时间性动态。
    \item 给出上下文驱动的任务分解不可靠的结构论证——这不是经验偶然（"经常出错"），而是信息论上的结构必然。
    \item 论证全自动Agent闭环的原则上的不可闭合性：层级验证链的根需求语义外在于模型的收敛空间，必须由系统外部提供。
\end{enumerate}


% ============================================================
\section{理论基础：注意力作为米田空间中的贝叶斯推断}

\subsection{分布语义学的范畴论表述}

分布语义学（distributional semantics）的核心直觉——一个词的意义由它与其他词的共现关系决定——在范畴论中有其自然的数学表达。

\begin{definition}[语义范畴]
    令 $\mathcal{C}$ 为语义范畴，其中：
    \begin{itemize}
        \item 对象 $\mathrm{Ob}(\mathcal{C})$ 为语义单元（概念、词义、或更一般的可被上下文区分的信息原子）；
        \item 态射 $\mathrm{Hom}_{\mathcal{C}}(A,B)$ 为从 $A$ 到 $B$ 的语义关系（上位词、同义、反义、共现、因果、时间顺序等）。
    \end{itemize}
\end{definition}

\textbf{米田引理（Yoneda Lemma）的语义学诠释。} 米田引理是范畴论中最基本的结构定理之一。在语义学背景下，它等价于以下陈述：

\begin{quote}
    一个语义单元 $X$ 由它与其他\textbf{所有}语义单元之间的关系\textbf{完全确定}。
\end{quote}

形式化地，对于任意 $X \in \mathcal{C}$，米田嵌入
\begin{equation}
    Y: \mathcal{C} \to [\mathcal{C}^{\mathrm{op}}, \mathbf{Set}], \quad
    Y(X) = \mathrm{Hom}_{\mathcal{C}}(-, X)
\end{equation}
将每个概念 $X$ 映射为"所有其他概念与 $X$ 的关系的集合"（即 $X$ 的"共现画像"的完整预层）。米田引理保证此嵌入是忠实全嵌入——即 $Y(X) \cong Y(Y)$ 当且仅当 $X \cong Y$。概念由其在全语义关系网络中的位置唯一确定。

\begin{definition}[语义空间]
    语义空间定义为米田嵌入的像所在的预层范畴 $[\mathcal{C}^{\mathrm{op}}, \mathbf{Set}]$——即所有概念在所有可能上下文中的关系画像所构成的空间。
\end{definition}

该空间在数学上是\textbf{无穷维}的——因为完备地刻画一个概念需要的态射（与其他概念的关系）数量原则上无界。这是§2.3截断投影概念的前置条件。

\subsection{注意力即贝叶斯推断}

Transformer的自注意力机制的核心计算为：
\begin{equation}
    \mathrm{Attention}(Q,K,V) = \mathrm{softmax}\!\left(\frac{QK^{\mathsf{T}}}{\sqrt{d_k}}\right)V
\end{equation}

其数学形式与Nadaraya-Watson非参数核估计同构：
\begin{equation}
    \hat{f}(x) = \sum_{i=1}^{n} \frac{K(x, x_i)}{\sum_{j=1}^{n} K(x, x_j)}\, y_i
\end{equation}

其中 $K(\cdot,\cdot)$ 为核函数，在注意力中对应 $\exp(QK^{\mathsf{T}}/\sqrt{d_k})$。

在高斯过程先验下，该估计器等价于后验均值的计算——即贝叶斯推断。因此，\textbf{单层单头注意力}在数学上等价于一次贝叶斯推断快照：以训练数据中习得的统计分布为隐式先验，对当前查询$Q$（提示词）条件下的最相关语义关联进行后验检索。

多层Transformer前馈传播则可以解释为多轮信念精炼（belief refinement）——每一层以上一层的输出为条件，在表示空间中更新语义关联的估计。其操作形式（前馈传播+残差连接）与标准贝叶斯条件更新的数学结构不完全等同，但精神相通：逐层减少语义不确定性。

这一类比有三个层次，各有不同的精度要求：
\begin{itemize}[leftmargin=*]
    \item \textbf{单层单头注意力} $\leftrightarrow$ 单次贝叶斯推断快照（精度较高：有Nadaraya-Watson同构支撑）；
    \item \textbf{多层Transformer前馈} $\leftrightarrow$ 多轮信念精炼（精度中等：残差传播与贝叶斯更新有结构类似但非等同）；
    \item \textbf{全程预训练} $\leftrightarrow$ 隐式贝叶斯模型平均（精度较低：是类比层面的统一）。
\end{itemize}

\subsection{米田嵌入的截断投影}

米田嵌入要求完备的态射集合——概念$X$的完整语义表征需要$X$与$\mathcal{C}$中\textbf{所有}对象的关系。在实际的Transformer中：

\begin{enumerate}[leftmargin=*]
    \item \textbf{参数空间的有效维度远小于名义维度。} Li et al.\ (2018)\ \cite{Li2018IntrinsicDim}通过子空间训练实验发现，神经网络目标景观的\textbf{本征维度}（intrinsic dimension）仅需数百到数千维即可达到良好性能，远小于模型的数百万甚至数十亿名义参数。换言之，参数空间的有效容量由本征维度而非名义维度决定。

    \item \textbf{低秩适应（LoRA）的有效性提供了独立证据。} Hu et al.\ (2021)\ \cite{Hu2021LoRA}证明，语言模型的权重更新矩阵可以通过极低秩（$r=1$--$16$）的分解来有效逼近。如果模型的完整参数空间中的有效信息可以用不到20维的子空间来承载增量变化，那么其表征的语义结构本身必然也是高度压缩的。

    \item \textbf{训练的本质是在低维流形上寻找最优投影。} Aghajanyan et al.\ (2020)\ \cite{Aghajanyan2020}的分析进一步确认了上述图景。梯度下降在数亿维的名义参数空间中搜索，但搜索轨迹的有效自由度仅限于本征维度决定的子流形。
\end{enumerate}

综合以上，\textbf{模型学到的语义映射不是完备的米田嵌入，而是米田嵌入在低维参数流形上的截断投影。} 从范畴论语言来看，存在一个因子分解：
\begin{equation}
    \xymatrix{
        \mathcal{C} \ar[r]^{Y} \ar[rd]_{\hat{Y}} & [\mathcal{C}^{\mathrm{op}}, \mathbf{Set}] \\
        & \mathcal{M}_d \ar[u]_{\iota}
    }
\end{equation}
其中 $\mathcal{M}_d$ 是维度为 $d$（本征维度）的参数流形，$\iota$ 是嵌入，$\hat{Y}$ 是模型实际学到的截断嵌入。忠实全嵌入的性质——概念由全部关系唯一确定——在截断后不再完整保持：两个在完整米田嵌入中可区分但相近的概念，可能在截断投影中重叠或混淆。

\begin{corollary}[核心约束]
    超出截断投影所覆盖的语义态射范围，模型无法提供结构保持的映射。这不是模型的"缺陷"——这是架构数学形式（softmax注意力 + 有限维参数空间）与语义空间的无限维结构之间的必然张力。
\end{corollary}


% ============================================================
\section{训练收敛：作为函子的学习过程}

\subsection{监督学习的函子表述}

将§2的范畴论框架置于动态视角下：训练过程的每一步，本质上是在构造一个保持关系结构的映射。

\begin{definition}[学习函子]
    令 $\mathcal{D}$ 为数据范畴（对象 = 带标注的数据点；态射 = 数据点之间的语义关系），$\mathcal{P}$ 为预测范畴（对象 = 模型输出；态射 = 输出之间的语义关系）。梯度下降训练的目标是找到一个映射
    \begin{equation}
        \mathcal{F}: \mathcal{D} \longrightarrow \mathcal{P}
    \end{equation}
    使得对于任意可复合态射 $f, g \in \mathrm{Mor}(\mathcal{D})$：
    \begin{equation}
        \mathcal{F}(g \circ f) \simeq \mathcal{F}(g) \circ \mathcal{F}(f)
    \end{equation}
    即 $\mathcal{F}$ 是一个\textbf{函子}——它保持数据中语义关系的复合结构。
\end{definition}

监督学习中的损失函数 $\mathcal{L}(\mathcal{F}(x), y)$ 衡量的正是该函子在每个训练样本上的结构保持程度。梯度下降沿损失景观的负梯度方向更新参数，在参数空间中搜索该函子的最优低维近似。

\subsection{收敛的边界}

监督损失驱动收敛到的是一个函数——该函数在\textbf{训练数据分布}上，以给定的损失度量，逼近最优的结构保持映射。但这个映射的全域范围由两重约束共同决定：

\begin{enumerate}[leftmargin=*]
    \item \textbf{分布约束：} 训练收敛仅在训练数据中\textbf{实际出现}的语义关系的分布范围内提供保证。训练数据中没有标注的语义关系——无论它们在客观的语义空间中是否成立——不在收敛保证的覆盖范围内。
    \item \textbf{截断约束（§2.3）：} 即使在训练分布内部，收敛到的映射也只是截断投影中的最优近似——模型表征的语义关系是完备米田嵌入的一个低维压缩。在这个压缩过程中丢失的语义信息，训练收敛无法恢复。
\end{enumerate}

\begin{proposition}[收敛边界]
    训练收敛的保证局限于截断投影覆盖范围内的训练数据语义分布。在此边界之外的语义态射，模型可能偶尔猜对，但其正确性不受任何训练收敛的数学保证。
\end{proposition}

这直接引出了一个关键区分：模型在训练分布内的局部态射保持可以可靠信赖；模型在分布外、或在需要复合多步态射的全局结构中的行为，不受同类保证。这是§5中训练收敛与测试收敛区分的数学基础。

\subsection{思维链RL的"补全"作用}

思维链强化学习（Chain-of-Thought RL）后训练能否扩展收敛的边界？我们的分析如下：

CoT RL后训练的实质是：在已有预训练模型的基础上，使用强化学习信号（过程奖励模型或结果验证）对推理路径进行优化。从范畴论视角来看：

\begin{enumerate}[leftmargin=*]
    \item \textbf{RL不扩展参数空间的维度。} 参数空间的本征维度由架构决定（§2.3），RL在同一个维度受限的参数流形上重新分布权重，而非增加流形的维度。
    \item \textbf{RL的作用是在截断投影\textbf{内部}重组映射结构。} 预训练阶段学到了一个截断投影 $\hat{Y}: \mathcal{C} \to \mathcal{M}_d$。该投影覆盖了某些语义态射，但这些态射的复合在该投影中的表述可能是不一致的。RL的作用是通过奖励信号显式地要求模型在多步推理中保持一致性——本质上是在 $\mathcal{M}_d$ 内部对被截断投影压缩的态射复合进行"拉直"（straightening）。
    \item \textbf{"补全"的是投影内部的一致性，而非投影本身。} CoT RL使原先隐式可推导但隐式推导不稳定（因为复合态射的一致性不受单独保证）的中间推理步骤变得显式化且可复现——但这发生在截断投影\textbf{之内}，而非扩展了投影的覆盖范围。RL无法"补全"整个米田嵌入中被截断投影丢失的维度。
\end{enumerate}

因此，CoT RL确实有效——它提升了模型在截断投影覆盖范围内的推理一致性——但它不能改变截断投影范围本身这一结构事实。这解释了为什么RL后训练的模型在分布内推理上大幅提升，但在需要超出训练语义分布的知识综合或创造性推理的任务上，提升有上限。


% ============================================================
\section{记忆的结构性二分}

\subsection{两种异质"记忆"}

LLM中存在两种常被统称为"记忆"但数学结构截然不同的信息存储机制。如表\ref{tab:memory_dichotomy}所示。

\begin{table}[htbp]
\centering
\caption{LLM记忆的结构性二分}
\label{tab:memory_dichotomy}
\begin{tabular}{@{}p{2.5cm}ll@{}}
\toprule
& \textbf{长期记忆} & \textbf{短期记忆} \\
\midrule
\textbf{载体}      & 模型参数（权重矩阵）      & 上下文窗口（token序列） \\
\textbf{形成方式}  & 训练收敛（一次性，静态）  & 顺序注入（每次推理重建） \\
\textbf{数学结构}  & 参数空间中的点（空间结构）& 序列中的线性展开（顺序结构） \\
\textbf{收敛保证}  & 有（训练损失驱动）        & 无（上下文映射不受收敛保证） \\
\textbf{动态性}    & 无（训练后冻结）          & 有顺序性，无时间性 \\
\textbf{更新方式}  & 离线重新训练/微调         & 在线追加/覆写 \\
\textbf{持久性}    & 跨会话持久                & 会话内挥发性 \\
\bottomrule
\end{tabular}
\end{table}

关键差异在于——这两种"记忆"不仅在体量或持久性上不同，它们的\textbf{数学结构}是不同的：

\begin{itemize}[leftmargin=*]
    \item \textbf{参数空间中的长期记忆}是一种\textbf{空间结构}：每个训练样本的贡献被压缩、混合并分布在整个权重矩阵中。各个事实不是以独立条目存储，而是以分布式表征的形式在参数空间中占据一个点。该点在训练后冻结——它代表的是训练数据的\textbf{全量统计压缩}。
    \item \textbf{上下文窗口中的短期记忆}是一种\textbf{顺序结构}：信息以线性序列的形式逐token注入，位置由位置编码标记。每次推理时从零开始——即使同一会话中，前一条消息也只能通过被追加到窗口中来"被记住"。
\end{itemize}

\subsection{无中间过渡的结构性论证}

2025年的多项工作（HAMR\cite{HAMR2025}、MemGPT、Memory Bear）提议在长期参数记忆和短期上下文记忆之间增设"中期记忆层"——通过摘要压缩、事件快照、向量检索等方式，构建一个跨越会话但非永久性的记忆层。

\textbf{我们的论证是：这一中期记忆层不存在——或者说，这些方案并未创建新的记忆\textbf{结构}，只是在管理短期记忆的\textbf{容量}。}

论证结构如下：

\begin{enumerate}[leftmargin=*]
    \item \textbf{大前提：} 长期记忆是空间结构，短期记忆是顺序结构——两者属于不同的数学范畴。空间结构中的信息以分布式表征的形式存在（"A点接近B点"），顺序结构中的信息以线性序列的形式存在（"A在B之前"）。
    \item \textbf{小前提：} 从顺序结构到空间结构，不存在连续的、可逆的映射。一个线性序列可以被压缩为一个固定长度的向量（例如摘要嵌入），但该压缩向量一旦形成，它就是一个\textbf{新的空间结构中的点}——它不再具有原序列的顺序可遍历性。反过来，一个空间结构中的点可以被"展开"为一段文本，但这段文本重新进入了顺序结构。
    \item \textbf{结论：} 不存在一个\textit{既具有空间结构的持久性又具有顺序结构的动态可遍历性}的中间记忆层。所谓的"中期记忆"方案本质上是两类操作：要么将顺序信息压缩进空间结构（摘要→嵌入→存储），此时它已成为静态的长期记忆载体（进入空间结构）；要么从存储中检索片段重新注入上下文窗口（检索→展开→注入），此时它被拉回了顺序结构。没有独立的"中间"数学结构。
\end{enumerate}

\textbf{澄清：我们不是在说这些工程方案"无效"。} 它们在各自的设定下有效——摘要压缩帮助容纳更多历史信息、向量检索帮助找到相关信息。但它们有效的方式是\textbf{管理短期记忆的容量}——即控制"什么信息进入有限的上下文窗口"——而非创建了一个新的、独立的记忆结构。

\subsection{时间性的缺失}

人类记忆具有丰富的时间性维度：Ebbinghaus遗忘曲线（随时间衰减）、睡眠巩固（时间依赖的离线重组）、情感调制（时间标记的情感强度影响提取概率）、干扰效应（新旧记忆竞争）、情境依赖提取（编码时的上下文影响回忆）。这些不是附加特征——它们是记忆之为\textbf{记忆}（而非存储）的构成性属性。

Transformer的位置编码提供的是\textbf{顺序}（$A$在$B$之前），不是\textbf{时间性}（$A$是三小时前说的，$B$是刚才说的）：

\begin{itemize}[leftmargin=*]
    \item 顺序是对称的：$A$在$B$前面和$B$在$A$后面是等价的，只是方向不同。
    \item 时间性是不对称的：三小时前的事和刚才的事在\textbf{提取的难易度、记忆的保真度、对当前行为的影响权重}上不同。刚才的事可能比三小时前的事更鲜活，但三小时前的事如果是关于一个关键决策，它的重要性可能更高——这两个维度不是线性可叠加的。
\end{itemize}

在Transformer中，所有位置的token在注意力机制下是平等可及的（通过并行的$QK^{\mathsf{T}}$计算），没有随时间衰减的提取难易度、没有离线巩固机制、没有自发的记忆重组。位置编码可以区分"A在B之前"和"B在A之前"，但无法编码"三个小时前"与"一分钟前"之间在\textbf{认知意义上的}差别。

\begin{corollary}[时间性缺失的工程后果]
    所有在现有Transformer架构内试图实现"渐进式持久记忆"的方案，面临一个共同的结构限制：它们可以更高效地管理上下文窗口的\textbf{容量}（哪些信息被保留、压缩和检索），但无法改变窗口内信息的\textbf{性质}——它仍然是顺序性的，没有时间厚度。这不是窗口大小的问题——即使把窗口扩到1亿token，信息仍然是按顺序排列的，没有随时间衰减的权重、没有巩固、没有重组。
\end{corollary}


% ============================================================
\section{测试收敛：Agent行为的真正支配原理}

\subsection{训练收敛不能解释Agent行为}

训练收敛（§3）描述的是：模型在参数空间中通过梯度下降找到最优低维函子近似的过程。这个过程在模型训练完成后\textbf{终止}——参数静态冻结。

Agent的行为与此不同：它不修改权重，而是在一个给定的上下文窗口内，通过多轮交互来逐步逼近目标。Agent做对（或做错）一件事的原因与训练收敛是\textbf{不同层面的因果关系}：

\begin{itemize}[leftmargin=*]
    \item 训练收敛解释为什么模型\textbf{有能力}做某类事——它在训练数据中学会了相关的局部语义态射。
    \item 训练收敛不解释为什么模型\textbf{在这一次}完成了（或偏离了）一个需要多步复合态射的复杂任务——因为复合后的一致性是全局结构属性，不由局部态射的保持性自动保证。
\end{itemize}

因此：训练得再好的模型，作为Agent时在多步任务中仍然会偏离。这不是训练不够好——这是驱动Agent行为的认知过程（测试收敛）与驱动模型能力形成的过程（训练收敛）具有不同的数学保证级别。

\subsection{测试收敛的定义与性质}

\begin{definition}[测试收敛]
    设Agent在一个上下文窗口内执行任务。令 $\mathcal{T}$ 为目标状态，令 $\mathcal{V}$ 为一组可自动执行的验证操作（满足条件→通过/不满足→未通过）。Agent通过循环的
    \[
    \text{生成潜在解} \longrightarrow \text{验证} \longrightarrow \text{分析差距} \longrightarrow \text{修正}
    \]
    逐步逼近 $\mathcal{T}$。这一过程的完成——即序列稳定在满足所有验证条件的状态——称为\textbf{测试收敛}。
\end{definition}

表\ref{tab:convergence_comparison}列出了训练收敛与测试收敛的系统性差异。

\begin{table}[htbp]
\centering
\caption{训练收敛与测试收敛的系统对比}
\label{tab:convergence_comparison}
\begin{tabular}{@{}p{3cm}ll@{}}
\toprule
& \textbf{训练收敛} & \textbf{测试收敛} \\
\midrule
\textbf{驱动信号}      & 监督损失（梯度）            & 可验证目标（测试通过/失败） \\
\textbf{收敛保证}      & 有（梯度下降的数学性质）     & 无（依赖验证信号的质量与覆盖度） \\
\textbf{作用范围}      & 参数空间（改变权重）         & 上下文空间（改变输出序列） \\
\textbf{持久性}        & 静态（训练后冻结）          & 挥发性（过程结束即丢弃） \\
\textbf{导向信号形式}  & 连续梯度向量场              & 离散的通过/失败信号 \\
\textbf{误差修正机制}  & 反向传播调整全部参数        & 前向修正下一轮输出 \\
\textbf{局部最优逃逸}  & 随机梯度噪声+SGD动量       & 无自动逃逸机制 \\
\bottomrule
\end{tabular}
\end{table}

\subsection{测试收敛的有效性条件}

测试收敛的有效性完全来自验证信号的覆盖度和精度。设任务目标 $\mathcal{T}$ 的约束集合为 $\mathcal{C} = \{c_1, c_2, \dots, c_n\}$，验证信号 $\mathcal{V}$ 覆盖了子集 $\mathcal{C}_{\mathcal{V}} \subseteq \mathcal{C}$。

\begin{proposition}[测试收敛的覆盖条件]
    当 $\mathcal{C}_{\mathcal{V}} = \mathcal{C}$（验证信号覆盖全部约束）时，测试收敛在原理上可以趋近 $\mathcal{T}$。当 $\mathcal{C}_{\mathcal{V}} \subsetneq \mathcal{C}$（存在未被验证信号覆盖的约束）时，Agent可以通过所有显式测试但最终输出不满足实际需求——这不是模型能力不足，而是验证信号的信息量不足以唯一确定目标状态。
\end{proposition}

工程实践中常见的"修一个bug引入三个新bug"正是这一机制的体现：验证信号覆盖了被修的那个bug，但没有覆盖被引入的三个新bug所在的约束维度。Agent通过了对当前验证集的测试收敛——它正确地完成了\textbf{它能被检验}的工作——但它偏离了它在当前验证信号下\textbf{看不见}的约束。

\subsection{与Test-time Scaling文献的关系}

2024--2025年间涌现的Test-time Scaling文献（SETS、IAD、ReVeal等）研究"如何在推理时多算几步以提升性能"。这些工作的核心思路是：将额外的计算量在推理时投入，通过搜索、采样、验证等方式提升输出质量。

\textbf{本文的追问与这些文献不同。} Test-time Scaling文献问："推理时多算几步\textbf{如何}提升性能？"本文追问道："推理时多算几步\textbf{为什么有时候有效有时无效}？"

本文的回答：有效与否取决于验证信号是否足够密集地覆盖目标结构的约束。在验证信号密集覆盖的任务上（如数学证明——每一步可自动验证），测试收敛的工作条件良好，test-time compute增加确实稳定提升性能。在验证信号稀疏或模糊的任务上（如"写一篇好文章"——"好"的语义约束无法被自动测试完全捕获），测试收敛的工作条件缺失，增加compute无法替代验证信号的消歧作用。


% ============================================================
\section{上下文边界与能力边界}

\subsection{文字输入的有限边界}

对于任何单次LLM推理（或一次连贯的多轮会话），模型能获取的\textbf{全部信息}来自两个来源：(1) 参数中的长期记忆（静态冻结）；(2) 上下文窗口中的短期记忆（顺序注入）。

上下文窗口定义了模型对单一问题的信息接入上限。但更关键的约束是：\textbf{有效利用长度远小于物理窗口容量}。Lost in the Middle效应\cite{Liu2024LostMiddle}已在多项研究中被系统性验证——窗口中间位置的信息检索准确率显著低于首尾位置。这种U形注意曲线是softmax注意力的结构特征（注意力权重通过softmax竞争归一化，首尾位置因边界效应获得更高权重），而非工程实现不优。

\begin{definition}[信息边界]
    令 $\mathrm{ContextCap}$ 为物理上下文窗口长度（以token计），$\mathrm{EffCap}$ 为模型能从上下文中\textbf{可靠提取结构保持映射}的最大信息量（以token计）。已知：
    \begin{equation}
        \mathrm{EffCap} \ll \mathrm{ContextCap}
    \end{equation}
    我们将 $\mathrm{EffCap}$ 称为\textbf{信息边界}——即单次任务中能被可靠处理的信息量上限。
\end{definition}

\subsection{边界外的任务必须分解}

如果一个任务所需的信息量超出信息边界，该任务在当前架构中无法被\textbf{单次连贯处理}——部分约束必然落在模型的有效注意范围之外。此时，任务必须被分解为更小的子任务，每个子任务的信息量控制在信息边界内。

\textbf{关键点：} 分解不是为了提高性能。分解是为了\textbf{恢复测试收敛的工作条件}——让每个子任务拥有独立的上下文窗口，在该窗口内信息量不超出边界，从而使测试收敛有运作的空间。

\subsection{分解本身也受边界约束}

分解行为——判断是否需要分解、如何分解、子任务间接口如何定义——本身也是一个需要信息处理的行为。分解行为发生在上下文窗口内，依赖上下文中的信息来做判断。但上下文中的信息不受收敛保证（§3.2、§4.1）——因此，\textbf{分解这一行为本身不受收敛保证}。

这构成了一个自指约束：

\begin{quote}
    判断一个任务是否需要分解，以及如何分解——这个判断所需的信息量，可能已经超出了上下文窗口的有效承载量。如果任务的信息量没有超出边界，就不需要分解；如果需要分解，分解决策所需的信息量本身可能已经超出了边界。
\end{quote}

这解释了以下经验观察的结构根源：长上下文模型（128K token甚至更长）在多步推理任务上的表现并未随窗口大小等比例提升——更大的窗口让模型"看到了更多"，但没有让模型"更可靠地处理了更多"。


% ============================================================
\section{任务分解的结构不可靠性}

\subsection{阶段节点向下一次性传导}

任务分解形成树形或图结构：父目标 → 子目标 → 孙目标 → 可执行原子任务。在此结构中：

\begin{proposition}[错误传导]
    设节点 $P$ 分解为子节点 $\{S_1, S_2, \dots, S_k\}$。若 $P$ 在语义层面存在偏差（$P$ 的指定不能唯一收敛到用户的真实意图），则所有子节点 $\{S_i\}$ 以该偏差为基础展开——偏差被继承和放大。子节点 $S_i$ 只持有 $P$ 赋予它的局部上下文，不具有独立检测和纠正 $P$ 层级错位所需的全域信息。
\end{proposition}

这类似于通信理论中的错误传播（error propagation），但在此处传播的不是比特错误而是语义错位——$P$ 的定义偏离了用户意图的"语义坐标"，$S_1$ 到 $S_k$ 在偏离后的坐标上各自展开，偏差沿层级累积而非抵消。

\subsection{上下文驱动分解不可靠的结构论证（三段论）}

\begin{enumerate}[leftmargin=*]
    \item \textbf{大前提：} 上下文环境中的信息不受收敛保证（§3.2、§4.1——上下文绕过训练收敛筛选，其中的信息没有经过参数空间的结构压缩）。
    \item \textbf{小前提：} 任务分解需要模型在上下文窗口内做出"统筹"和"连贯"的判断——即需要可靠的复合态射保持。模型需要识别全局依赖结构、判断哪些子任务耦合、在何处划分接口——这些判断如果需要被信任，其背后的推理必须保持语义关系在复合后的一致性。
    \item \textbf{结论：} 使用不受收敛保证的信息来做出需要收敛保证的决策——结构不可靠是逻辑上的必然，而非经验上的偶然。
\end{enumerate}

\begin{corollary}
    在Agent系统的设计中，将任务分解的\textbf{结构决策}（架构选择、接口定义、语义分歧消解）交由模型在上下文中自主完成——这种设计的失败率有理论下界。它不是"prompt写得不够好"——它在当前架构的数学结构下不可能被写得足够好。
\end{corollary}

\subsection{与经验文献的归因层差异}

2025年的多项经验研究从现象层面验证了任务分解的脆弱性：

\begin{itemize}[leftmargin=*]
    \item LEAD\cite{LEAD2025}识别了多步规划中的系统性偏离，归因为规划-执行脱节；
    \item Divide-and-Conquer噪声模型将分解失败归因为子任务边界的噪声积累；
    \item CoDA报告了"上下文爆炸"——分解产生的子上下文总量超出管理能力。
\end{itemize}

\textbf{本文的归因与上述文献不同。} 它们将脆弱性归因为噪声、容量、协调机制——这些是"如果改进X，问题可以缓解"的归因。本文将脆弱性归因为\textbf{上下文本身不受收敛保证}——这是"改进X不能解决根本问题，因为X没有触及因果链的根节点"的归因。

这不是说上述文献的归因错误——噪声积累和上下文过载确实是现象层面的机制。但噪声归因暗示"去噪即可"，结构归因指出：即使没有噪声，用无保证的信息来做需保证的决策，逻辑上就是在掷骰子。噪声降低可以提升掷骰子的方差（使失败不那么离谱），不能改变掷骰子这一根本性质。

\subsection{关键分叉点须外部介入}

从上述论证的直接推论是：在任务分解的\textbf{不可逆决策节点}上，模型在上下文中的自主判断不可信任。\textbf{不可逆}在此处的精确含义是：决策一旦做出并被下传到子节点，子节点不具备独立的全局信息来纠正该决策。

这些关键分叉点包括但不限于：
\begin{itemize}[leftmargin=*]
    \item 对用户需求的模糊表述进行消歧和具体化（两种消歧方向导致完全不同的任务树）；
    \item 顶层子任务的接口定义（父层接口偏差使所有子层偏离）；
    \item 子任务间的依赖关系声明（遗漏耦合导致下层集成失败）。
\end{itemize}

在这些节点上，必须由\textbf{系统外部}——人工操作者，或至少一个独立于模型上下文窗口的硬约束程序（如类型检查器、编译器、自动化测试套件）——进行仲裁。工程上，这对应了当前最成功的Agent系统（Cursor、Claude Code）的设计模式：在关键决策点设置人工确认，全自动早期Agent（AutoGPT）几乎全部失败。


% ============================================================
\section{层级收敛与覆盖面的指数扩张}

\subsection{分解的层级结构}

给定一个超边界任务，Agent（在外部仲裁的参与下）将其分解为层级需求结构：

\[
\text{根需求（用户给定）} \longrightarrow \text{阶段需求} \longrightarrow \text{子阶段需求} \longrightarrow \text{可验证原子单元}
\]

每一层级的分解行为本身是在\textbf{父级上下文窗口}中完成的。制定完成后，各子需求释放到各自的独立上下文窗口中，各自完成其局部的测试收敛。由此形成多级收敛的嵌套结构。

\subsection{覆盖面的指数扩张与收敛代价}

层级结构的工程价值在于：每个节点可以独立地收敛到其局部可验证目标，而不需要在全局上下文中同时维护所有约束。覆盖面随层级指数增长——$k$级分解产生$O(b^k)$个可并行的验证单元（$b$为平均分支因子）。

但代价也是结构性的：

\begin{itemize}[leftmargin=*]
    \item \textbf{高层收敛成本仍受窗口约束。} 高层级（根需求→阶段需求的分解）的上下文窗口需要容纳全局结构——子任务划分、接口定义、依赖关系——随着任务复杂度增加，高层上下文同样会触及信息边界。
    \item \textbf{验证标准沿层递归上溯。} 每一级只看到其子级的\textbf{验证结果}（通过/失败），而非子级的完整上下文。这意味着上层对下层的信息损失是累积的——$n$级层级结构引入$n$层信息压缩。
\end{itemize}

\subsection{收敛的终止点：根需求}

验证标准沿层级链向上追溯，最终到达\textbf{根需求}——"此结果是否符合用户的实际意图"。这一判断的语义坐标——具体用户在具体时刻的具体需求——不在模型的收敛空间内。原因如下：

\begin{enumerate}[leftmargin=*]
    \item 模型的收敛空间由训练数据分布定义（§3.2）；
    \item "具体用户此刻的具体意图"是一个\textbf{单点语义事件}——它不出现在任何训练数据中；
    \item 因此它在米田嵌入中的表征不在截断投影的覆盖范围内（§2.3）。
\end{enumerate}

\begin{proposition}[根需求的外在性]
    层级验证链的终止点——根需求的真值条件——的语义\textbf{原则上}外在于模型的收敛空间。这不是模型不够大、训练数据不够多的问题——这是收敛空间的结构边界决定的。根需求必须由\textbf{系统外部}（用户、物理世界、或独立于模型的硬约束源）提供。
\end{proposition}


% ============================================================
\section{全自动闭环的原则不可闭合性}

\subsection{不可闭合的精确含义}

"全自动闭环在原则上不可闭合"这一断言需要精确限定范围——不精确的表述会超出本文实际论证的承载能力。三层区分如下：

\begin{table}[htbp]
\centering
\caption{三层不可闭合性}
\label{tab:three_layers}
\begin{tabular}{@{}p{2.5cm}p{4cm}p{2.5cm}p{2.5cm}@{}}
\toprule
\textbf{层面} & \textbf{定义} & \textbf{不可闭合的性质} & \textbf{论证强度} \\
\midrule
\textbf{目标设定层} & 系统能否自主生成最高目标 & 逻辑必然不可自举 & \textbf{强} \\
& & 目标语义不在训练分布内，& \\
& & 无从收敛 & \\
\midrule
\textbf{任务分解层} & 系统能否在给定目标下自主规划子任务 & 结构不可靠 & \textbf{中} \\
& & 上下文驱动的分解决策 & \\
& & 不受收敛保证（§7） & \\
\midrule
\textbf{执行闭环层} & 系统能否在给定明确验证标准下自主完成任务 & 条件可行 & \textbf{弱} \\
& & 若验证信号覆盖充分，& \\
& & 执行可在无人工介入下完成 & \\
\bottomrule
\end{tabular}
\end{table}

\textbf{本文的核心论证支持目标设定层的不可闭合性。} 执行闭环层在条件满足时可以在无需人工介入的情况下完成——例如，给定完整的自动化测试套件和明确的接口规约，Agent可以独立完成代码实现并通过所有测试。这不与该论证矛盾——它论证的是"最高目标不可自举"，而非"任何执行步骤都不能自动化"。

\subsection{与现有AGI不可能论证的比较}

已有文献从其他角度论证了通用AI或全自动Agent的原则限制：

\begin{itemize}[leftmargin=*]
    \item \textbf{哥德尔路径：} 任何形式系统存在不可证的真理（G\"odel, 1931）$\to$ AI受限于其所运行的形式系统。
    \item \textbf{图灵路径：} 停机问题不可判定（Turing, 1936）$\to$ 存在AI无法提前判定的程序行为。
    \item \textbf{中文屋路径：} 符号操作不等于理解（Searle, 1980）$\to$ 纯形式系统不具备意向性。
    \item \textbf{热力学/信息论路径：} Landauer原理设定计算能量下限；信息不可无代价获取。
\end{itemize}

\textbf{本文的路径与以上均不同。} 本文从以下事实出发：
\[
\text{收敛空间的有界性（§3.2）} \longrightarrow \text{米田嵌入的截断性（§2.3）} \longrightarrow \text{根需求外在于截断投影（§8.4）}
\]
这一路径的优势在于：它不依赖"智能"或"理解"的哲学定义，不依赖计算复杂度的极限假设，只依赖一个经验前提（参数空间有效维度有限）和一个数学类比（注意力与贝叶斯推断的同构），推出一个关于Agent闭环的约束性结论。如果前提被新的实验证据修正（例如发现参数空间的本征维度可以动态无限扩展），结论也需要相应修正——这使该论证在原则上\textbf{可证伪}。


% ============================================================
\section{Agent = 模型 + Harness：工程原则}

\subsection{定义与理论角色}

\begin{definition}[Agent]
    $\mathrm{Agent} = \mathrm{Model} + \mathrm{Harness}$，其中：
    \begin{itemize}
        \item $\mathrm{Model}$（模型）提供截断的语义映射能力（训练收敛的产物）；
        \item $\mathrm{Harness}$（驭具）提供行动框架：$\mathrm{Harness} = \mathrm{程序框架} + \mathrm{上下文环境}$。
    \end{itemize}
    模型与Harness的关系类似于远程操控者与机甲的关系——模型提供认知能力，Harness提供外部验证、硬约束和状态管理。两者不属于同一个收敛系统：模型的能力由训练收敛决定，Harness的行为由测试收敛（和/或程序化确定性规则）驱动。
\end{definition}

Harness不是"模型能力不够时的权宜之计"——它是\textbf{收敛理论在工程层面的实现}。在模型的结构边界不会因能力提升而消失的前提下（因为边界来自于架构数学结构而非训练不充分），Harness的设计是第一性的，不是过渡性的。

\subsection{五项根本工程原则}

基于前述完整推理链（§2--§9），我们导出Agent工程设计的五项根本原则：

\begin{enumerate}[leftmargin=*,label=\textbf{P\arabic*}.]
    \item \textbf{精确性优先（精）：} 提示词的明确指示——最大限度消除输入在语义空间中的歧义维度。准确有限维的输入更容易找到被收敛保证的结构保持映射。\textbf{注：} "精确"不等于"短"——CoT增加了token数但通过消歧降低了信息结构复杂度，在此意义上是更精确的（参见§5.4中的"表面复杂度"与"结构复杂度"区分）。

    \item \textbf{结构简化（简）：} 最小化输入的信息结构复杂度——即最小化语义歧义维度的数量，而非最小化token数。每增加一个可被多重解释的信息维度，输入在语义空间中的可能落点数量指数增长，超出截断投影覆盖范围的概率随之增加。

    \item \textbf{语义氛围（氛围）：} 将输入置于合适的语义氛围——即选择与目标态射最近的训练分布区域，利用训练收敛在该区域内的高密度覆盖来降低被截断投影遗漏的概率。

    \item \textbf{分叉控制：} 关键决策节点（任务划分、接口定义、需求消歧）必须由外部仲裁介入——不信任上下文驱动的结构决策（§7论证）。

    \item \textbf{可验证性：} 分解后的每个任务单元必须配备可自动执行的验证标准——因为Agent行为由测试收敛驱动，而测试收敛依赖验证信号的存在和覆盖度（§5论证）。
\end{enumerate}

\subsection{可行域边界}

上述原则共同划定了当前Transformer架构下Agent工程设计的\textbf{可行域}：

\begin{itemize}[leftmargin=*]
    \item \textbf{可行：} 将任务切至上下文信息边界内，使测试收敛有条件运作；每个单元配备可自动执行的验证标准；结构决策节点由外部仲裁介入；接受根需求的外部性——人机协作在此结构下不是权宜的过渡方案，而是结构的必然要求。
    \item \textbf{不可行（在当前架构的数学约束下）：} 依赖模型自主规划长程任务而不设外部仲裁点；以模糊需求驱动多步自主决策而不设自动验证；试图在现有Transformer架构内实现完整的闭环自举。
\end{itemize}

\textbf{关键澄清：} "不可行"在此语境下不是"技术上困难"——它是"给定架构的数学结构，该目标的收敛条件在原则上无法满足"。如果未来的架构创新改变了前提（例如参数空间的有效维度可以动态扩展至覆盖根需求语义），上述不可行性边界需要相应重划。


% ============================================================
\section{相关工作}

本节逐一标注本文与已有文献的差异化路径。差异化不等于否定——已有文献的经验发现本文大多接受，我们的差异在于归因的深度和统一性。

\subsection{范畴论与深度学习}

范畴论应用于神经网络/深度学习的文献包括：
\begin{itemize}[leftmargin=*]
    \item \textbf{Self-Attention as a Parametric Endofunctor} (O'Neill, 2025)：将自注意力构造为参数化自函子，使用自由单子（free monad）框架。与本文的差异：(1) 本文使用米田嵌入而非自由单子作为核心构造；(2) 本文从范畴结构\textbf{推导工程约束}，而非仅停留在数学表述。
    \item \textbf{Categorical Invariants of Learning Dynamics} (Tamim, 2025)\ \cite{Tamim2025Categorical}：将学习定义为函子 $\mathcal{L}: \mathbf{Param} \to \mathbf{Rep}$，运用同伦论分析训练轨迹。与本文最接近的独立工作，但：(1) 未涉及米田嵌入；(2) 未将函子概念延伸至Agent行为层面。
\end{itemize}

\subsection{Agent任务分解}

\begin{itemize}[leftmargin=*]
    \item \textbf{LEAD} (Pushkin \& Abbe, 2025)\cite{LEAD2025}、\textbf{Divide-and-Conquer噪声框架}、\textbf{CoDA}等：从经验层面识别和分析了多步任务分解中的失败模式，归因为噪声积累、规划-执行脱节、上下文过载。本文接受这些经验发现，但论证这些失败的根本原因比噪声更深一层——在于上下文信息\textbf{本身不受收敛保证}。
    \item \textbf{Chain-of-Agents} (NeurIPS 2024)：多Agent协作框架，将复杂度从$O(n^2)$降至$O(nk)$。与本文§8的层级收敛框架在精神上一致，但CoA未从其架构有效性中提炼"测试收敛"这一一般原理。
\end{itemize}

\subsection{LLM记忆}

\begin{itemize}[leftmargin=*]
    \item \textbf{Cognitive Memory in LLMs} (2025)\cite{CognitiveMemory2025}、\textbf{HAMR}\cite{HAMR2025}、\textbf{Memory Bear}等：在工程架构层面提出多级记忆方案（工作记忆/短期/中期/长期）。本文的差异化贡献：从数学结构的异质性出发，论证\textbf{不存在独立的"中期记忆"结构}——这些方案有效是因为它们管理了短期记忆的容量，而非创建了新的记忆结构。
    \item \textbf{Titans架构} (Meta, 2025)：提出在推理时使用"惊喜"梯度更新神经记忆模块。从本文框架的角度，Titans试图在顺序结构中注入有限的空间结构更新——但该更新是推理时的一次性前馈操作，不具有在线巩固、重组、时间衰减等人类记忆的时间性特征。
\end{itemize}

\subsection{Test-time Scaling}

SETS、IAD、ReVeal等文献研究推理时的计算扩展策略。本文的差异化：不是研究如何扩展推理时计算，而是分析\textbf{推理时计算为什么有效}——提出测试收敛作为统一的机理解释，并分析其有保证性缺失这一根本限制。

\subsection{AGI理论局限}

G\"odel/Turing/Searle/热力学路径各有其理论前提和限定。本文的路径从实验上可检验的前提（参数空间有效维度有限）出发，推出模型收敛空间的结构边界——\textbf{论证的起点不是哲学假设而是可测量的经验量}。如果未来的实验发现参数空间的本征维度可随模型规模以某种方式无限增长，本文的部分结论需要修正——这一可证伪性本身就是本文与不可证伪的哲学论证的关键区别。

\subsection{参数空间本征维度}

Li et al.\ (2018)\cite{Li2018IntrinsicDim}、Aghajanyan et al.\ (2020)\cite{Aghajanyan2020}、LoRA\cite{Hu2021LoRA}——这些工作为本文的"截断投影"命题提供了独立于人文本框架的经验基础。这些工作本身未将其发现与LLM Agent的行为边界或语义空间的范畴结构建立联系——这是本文的独有综合。


% ============================================================
\section{讨论与展望}

\subsection{本文论证的可检验推论}

如果本文框架为真，应观察到以下经验模式。这些推论附带了\textbf{证伪条件}——它们的可检验性是本文框架可证伪性的具体体现。

\begin{enumerate}[leftmargin=*,label=\textbf{T\arabic*}.]
    \item \textbf{上下文窗口扩展不会消除任务分解的结构不可靠性。}
    预测：将上下文窗口从128K扩展至1M乃至10M token，多步长程规划任务的自主分解失败率不会降为零，而会收敛到某个非零下界。该下界来自于信息边界的结构约束，而非容量约束。
    证伪条件：如果观察到随窗口扩展，多步规划自主分解失败率持续下降且无收敛至非零下界的迹象，则§6--§7的论证需要修正。

    \item \textbf{现有架构内的记忆增强方案只能改进容量管理，不能创造时间性动态。}
    预测：任何在Transformer架构内（不改变底层的softmax注意力+参数静态冻结）的"中期记忆"方案，在长时间跨度的多会话任务中，其性能增益的来源可以被完全归因为更优的检索和摘要（即容量管理），而非真正的时间性动态（遗忘曲线、离线巩固、自发重组）。
    证伪条件：如果发现一个方案在现有架构内实现了具有Ebbinghaus遗忘曲线特征的、随时间自发衰减和重组的记忆（而非通过prompt显式注入"时间流逝"信息），则§4的论证需要修正。

    \item \textbf{全自动Agent在开放域中的闭环成功率存在系统性上界。}
    预测：该上界由验证信号的覆盖度决定。在某类开放域任务上（如：从模糊自然语言需求出发的端到端软件开发），即使模型能力持续提升（参数规模、训练数据量、RL后训练），全自动闭环成功率的上界将显著低于半自动（人工在关键分叉点介入）方案。
    证伪条件：如果未来出现一个在模糊自然语言需求下可靠完成端到端全自动软件开发的Agent（成功率与带人工确认的版本无显著差异），则§9的论证需要修正。
\end{enumerate}

\subsection{对架构设计的启示}

如果本文的框架成立，它对下一代架构设计有以下启示：

\begin{enumerate}[leftmargin=*]
    \item \textbf{如果目标是时间性记忆，需要从根本上改变记忆的数学结构。} 在softmax注意力中引入显式的时间衰减权重、离线巩固阶段（类似于睡眠期间的记忆重组）、或竞争性提取机制（新信息与旧信息的干扰效应）——不只是加一个"记忆模块"。

    \item \textbf{Harness的设计是第一性的，不是过渡性的。} 因为模型的结构边界不会因模型能力提升而消失（边界来自架构数学结构，而非训练不充分），Harness不是"当前模型还不够好时的权宜之计"。外部仲裁、可验证单元、层级分解——这些是在既有架构下Agent可靠运行的构成性条件，而非可被后续模型升级淘汰的临时方案。

    \item \textbf{人机协作不是过渡形态。} 在目标设定的不可自举性被架构创新改变之前（即参数收敛空间能动态覆盖根需求语义），人在任务设定和关键分叉决策中的角色是结构性的——它可以被优化（更高效的交互、更精准的确认点），但不可以被消除。
\end{enumerate}

\subsection{局限与开放问题}

本文的诚实标注：

\begin{enumerate}[leftmargin=*]
    \item \textbf{米田嵌入与贝叶斯推断的数学对偶。} §2中二者的统一目前处于类比层面——两者在结构上的对应已被识别，但尚未给出严格的数学构造。开放问题：是否可能存在一个交换图，将注意力的softmax计算、Nadaraya-Watson核估计与米田嵌入统一为同一个数学结构的不同投影？这是本文框架中最需要后续数学工作的部分。

    \item \textbf{测试收敛的精确信息论刻画。} §5提出了测试收敛的"保证性缺失"，但未将其量化为精确的信息论条件。开放问题：给定目标约束集合 $\mathcal{C}$ 和验证信号 $\mathcal{V}$，测试收敛成功的概率是否有关于 $|\mathcal{C} \setminus \mathcal{C}_{\mathcal{V}}|$ 的精确函数形式？验证信号的最小信息量是多少？

    \item \textbf{本征维度的可扩展性。} §2.3和§3.2的核心前提——参数空间有效维度有限——是一个经验事实而非逻辑必然（本征维度可能随模型规模扩展而增长，虽然当前证据表明增长率远低于名义参数增长率）。如果未来的架构（如动态稀疏激活、持续学习）使有效维度可以持续扩展至接近或超越根需求语义的复杂度，则本文基于截断投影的收敛边界论证需要修正。

    \item \textbf{本文不讨论意识、理解或智能的哲学定义。} 本文的全部论证聚焦于：在给定架构的数学结构下，模型能可靠执行的操作集合的边界。这些边界是否意味着模型"没有智能"或"不理解"——不在本文的论证范围内。本文的策略是：在哲学问题被解决之前，先回答一个更窄但可操作的问题——\textbf{边界在哪里，为什么在那里}。
\end{enumerate}


% ============================================================
\section{结论}

本文提出了一个从Transformer数学结构到Agent工程约束的统一收敛理论框架。核心论证链为：

\[
\boxed{
\begin{array}{c}
\text{注意力 = 语义空间中的贝叶斯推断（Nadaraya-Watson同构）} \\
\Downarrow \\
\text{参数空间有效维度有限（Li et al. 2018; LoRA）} \Longrightarrow \text{截断投影} \\
\Downarrow \\
\text{训练收敛仅覆盖截断投影内的语义态射} \\
\Downarrow \\
\text{上下文中的信息绕过训练收敛筛选} \Longrightarrow \text{不受收敛保证} \\
\Downarrow \\
\text{Agent行为 = 上下文中的测试收敛} \Longrightarrow \text{缺乏数学保证} \\
\Downarrow \\
\text{分解决策依赖上下文} \Longrightarrow \text{结构不可靠（三段论）} \\
\Downarrow \\
\text{验证链上溯至根需求} \Longrightarrow \text{根需求外在于收敛空间} \Longrightarrow \text{不可自举}
\end{array}
}
\]

每一步的上层前提标于箭头上方。如果某一步的前提在未来的理论或实验工作中被修正，结论也需要相应修正——这一可修正性本身就是本文框架与不可证伪的哲学断言的关键区别。

\textbf{本文的中心主张：} LLM Agent面临的若干核心困难——任务分解的结构性脆弱、记忆缺乏时间深度、全自动闭环的不可达——不是工程上的暂时问题。它们是Transformer架构的数学结构（有限维参数空间上的截断语义投影）在Agent多步行为中展开时的必然结果。理解这些边界的来源和性质，不是限制Agent工程的雄心，而是使其建立在一个经得起现实检验的基础上。

\vspace{12pt}

\noindent\textbf{致谢。} 本文的初稿结构受益于与多个大型语言模型的对话，这些对话中的批判性反驳迫使作者逐条明确各论证的边界、强度和薄弱环节。所有错误和未完成的论证由作者负责。

% ============================================================
% 参考文献
% ============================================================

\begin{thebibliography}{99}

\bibitem{Vaswani2017}
A.~Vaswani et al.\ (2017).
``Attention Is All You Need.''
\textit{NeurIPS 2017}.

\bibitem{Li2018IntrinsicDim}
C.~Li, H.~Farkhoor, R.~Liu, J.~Yosinski (2018).
``Measuring the Intrinsic Dimension of Objective Landscapes.''
\textit{ICLR 2018}. arXiv:1804.08838.

\bibitem{Aghajanyan2020}
A.~Aghajanyan, L.~Zettlemoyer, S.~Gupta (2020).
``Intrinsic Dimensionality Explains the Effectiveness of Language Model Fine-Tuning.''
\textit{ACL 2021}. arXiv:2012.13255.

\bibitem{Hu2021LoRA}
E.~J.~Hu et al.\ (2021).
``LoRA: Low-Rank Adaptation of Large Language Models.''
\textit{ICLR 2022}. arXiv:2106.09685.

\bibitem{Liu2024LostMiddle}
N.~F.~Liu et al.\ (2024).
``Lost in the Middle: How Language Models Use Long Contexts.''
\textit{TACL 2024}. arXiv:2307.03172.

\bibitem{Tamim2025Categorical}
K.~Tamim et al.\ (2025).
``Categorical Invariants of Learning Dynamics.''
arXiv:2510.04376.

\bibitem{LEAD2025}
S.~Pushkin, E.~Abbe (2025).
``LEAD: Language-based Error Analysis for Decomposition in LLM Agents.''
\textit{Preprint}, 2025.

\bibitem{HAMR2025}
Z.~Adams (2025).
``HAMR: Hierarchical Adaptive Memory Retrieval for LLM Agents.''
\textit{Preprint}, 2025.

\bibitem{CognitiveMemory2025}
Y.~Wu et al.\ (2025).
``Cognitive Memory in Large Language Models.''
arXiv:2504.02441.

\bibitem{MACI2025}
E.~Y.~Chang et al.\ (2020--2025).
``MACI: Multi-Agent Collaborative Intelligence Framework.''
Stanford University.

\bibitem{ChainOfAgents2024}
Y.~Zhang et al.\ (2024).
``Chain of Agents: Large Language Models Collaborating on Long-Context Tasks.''
\textit{NeurIPS 2024}.

\bibitem{Anthropic2024Agents}
Anthropic (2024).
``Building Effective Agents.''
Internal technical report.

\bibitem{SignificanceDeficit2025}
A.~Shillington et al.\ (2025).
``The Significance Deficit: A Unified Theory of Transformer Failures.''
Zenodo. 10.5281/zenodo.17847870.

\bibitem{Kaiser2025FirstPrinciples}
L.~Kaiser (2025).
``大模型的第一性思考.'' 36Kr / 公开演讲.

\bibitem{KarpathyContextRAM}
A.~Karpathy (2024).
``LLM OS: Context is RAM.''
Public writing.

\bibitem{Xu2026Holographic}
Y.~Xu (2026).
``Global Low-Rank, Local Full-Rank: The Holographic Encoding of Learned Algorithms.''
arXiv:2602.18649.

\bibitem{CoA_NeurIPS2024}
``Chain-of-Agents'' (2024).
\textit{NeurIPS 2024}. arXiv:2406.02818.

\end{thebibliography}

\end{document}
